%% sections/intro.tex
%% Data-centric (DMLR) framing: foreground value to ML research, not just SE tooling.

\section{Introduction}
\label{sec:intro}

Source code is one of the most heavily used data modalities in machine learning, and
the dominant way to obtain it is to mine it: scrape repositories, filter, deduplicate,
and ship. Mined corpora are abundant, but they are \emph{uncontrolled}. A snippet's
semantics, its surface ``messiness,'' and its difficulty are whatever the wild happened
to contain; there is no known-optimal reference to grade an instance against; and the
two things a study most often wants to vary --- \emph{how hard the problem is} and
\emph{how poorly it is written} --- are confounded, because in real code a longer
program is usually also a messier one. Worse, every public sample is a moving target for
contamination: any snippet a researcher uses to evaluate a code model may already sit in
that model's training set, so a high score can reflect memorization rather than
capability~\citep{contamination-placeholder}. Mining gives us volume; it does not give us
\emph{control}, \emph{labels}, or \emph{freshness}.

A growing set of data-centric ML problems need exactly that control. Contamination-aware
evaluation of code models requires items that can be \emph{re-minted} after a training
cutoff rather than drawn from a static, leakable pool. Training and probing studies want
parallel data in which one factor moves while the rest are held fixed. Robustness
analysis wants to ask whether a model degrades because \emph{the problem got bigger} or
because \emph{the presentation got messier} --- a question no fixed corpus can answer,
because the two grow together in the wild. Refactoring, clone-detection, and
code-translation research wants paired clean/messy programs whose semantic equivalence is
\emph{guaranteed}, not merely assumed from a commit message. Each of these is a request
for a piece of data that mined corpora structurally cannot supply: a labelled,
equivalence-controlled, freshly-mintable, multi-language instance.

We present \tool, a tool that synthesizes such data. \tool{} is an
\emph{anti-optimization transpiler}: where an optimizing compiler holds semantics fixed
and \emph{minimizes} a cost, \tool{} holds semantics fixed and deliberately
\emph{maximizes} incidental complexity. It maps a clean, language-agnostic JSON
intermediate representation (IR) to deliberately redundant, fully-flattened programs in
five languages (Python, JavaScript, Go, Java, C++), and every emitted program is
compiled, run, and checked against a reference oracle, so each instance is correct
\emph{by construction}. Because the clean IR is a known-optimal reference and
``messiness'' is dialed by composable, strictly-nested anti-pattern profiles, each
instance carries labels along two \emph{orthogonal} difficulty axes --- \emph{intrinsic}
(the problem genuinely grows) and \emph{incidental} (presentation degrades at fixed
semantics) --- and a single IR yields a parallel corpus across all five languages. To
resist contamination, scored instances are minted fresh from a \emph{private} held-out
seed that is never committed, a public canary GUID enables training-set audits, and
graded novelty tiers turn ``how memorized is this?'' into a measured quantity rather than
a hope. A motivating downstream use case --- a meta-evaluation of whether LLM judges
track a code-quality order known independently of any rater --- is the subject of a
separate, deferred study~\citep{TODO-companion-study}; here we deliver and document the
\emph{instrument} and the dataset it produces.

\paragraph{Contributions.} This paper contributes:
\begin{itemize}
  \item \textbf{The generator and its correctness-by-construction design}
        (Section~\Cref{sec:resource}): a staged IR$\rightarrow$five-language pipeline in which
        valid syntax is structural and runtime safety is always-on, gated by an oracle
        validator so that every released instance provably builds and computes the
        intended result.
  \item \textbf{Two orthogonal, by-construction difficulty labels}
        (Section~\Cref{sec:axes}): an \emph{incidental} messiness knob (a strictly-nested
        anti-pattern ladder) crossed with an \emph{intrinsic} problem-size knob, so a
        score change can be attributed to one cause with the other held fixed --- a
        decomposition no mined corpus offers.
  \item \textbf{A contamination protocol} (Section~\Cref{sec:contam}): mint-after-cutoff
        generation, a public \texttt{dev} split versus a private held-out \texttt{test}
        split, an embedded canary, and graded A/B/C novelty tiers --- with the
        public/private asymmetry as the mitigation, not an apology.
  \item \textbf{Construct-validity evidence} (Section~\Cref{sec:validity}) that the
        by-construction quality order moves established complexity and readability
        metrics, reported with its weaknesses stated plainly.
  \item \textbf{Reference baselines} (Section~\Cref{sec:baselines}) that establish the
        dataset is usable and discriminative, and a reusable, FAIR, DOI-archived,
        dependency-free open-source artifact.
\end{itemize}

\paragraph{Artifact.} \tool{} is open source under the MIT license, depends only on the
Python standard library, ships runnable examples, a test suite, and a documented
extension contract, and is archived with a DOI; a ``Datasheet for Datasets''
(Appendix~\Cref{app:datasheet}) accompanies the dataset.
